\exercisesection{2}

\begin{enumerate}
\def\labelenumi{\arabic{enumi}.}
\item
  Give an example of an exact sequence \(0 \rightarrow N' \rightarrow N \rightarrow N'' \rightarrow 0\) of left A-modules and a right A-module E such that E is \(N'\) -flat and \(N''\) -flat, but not N-flat (take for example \(N' = N'' = Z/2Z\) ).
\item
  Let M, N be two submodules of an A-module E such that \(M + N\) is flat. For M and N to be flat, it is necessary and sufficient that \(M \cap N\) be flat.
\item
  Let A be the ring K{[}X, Y{]} of polynomials in two indeterminates over a field K.
\end{enumerate}

\begin{enumerate}
\def\labelenumi{(\alph{enumi})}
\item
  Consider in A the principal ideals \(\mathfrak{b} = (\mathbf{X}), \mathfrak{c} = (\mathbf{Y})\) , which are free A-modules and whose intersection \(\mathfrak{b} \cap \mathfrak{c} = (\mathbf{XY})\) is also free. Show that \(\mathfrak{a} = \mathfrak{b} + \mathfrak{c}\) is not a flat A-module, although \(\mathfrak{a}\) is torsion-free (cf.~Algebra, Chapter III, § 2, Exercise 4).
\item
  In the A-module \(A^{2}\) , let R be the submodule consisting of the elements \((x, -x)\) where \(x \in a\) . In the A-module \(A^{2}/R\) , let M, N be the submodules which are the images of the factor submodules of \(A^{2}\) ; show that M and N are isomorphic to A, but that \(M \cap N\) is not a flat A-module.
\end{enumerate}

\begin{enumerate}
\def\labelenumi{\arabic{enumi}.}
\setcounter{enumi}{3}
\tightlist
\item
  \begin{enumerate}
  \def\labelenumii{(\alph{enumii})}
  \tightlist
  \item
    Give an example of an exact sequence
  \end{enumerate}
\end{enumerate}

\[
0 \rightarrow \mathrm{E} ^ {\prime} \rightarrow \mathrm{E} \rightarrow \mathrm{E} ^ {\prime \prime} \rightarrow 0
\]

which does not split and all of whose terms are flat modules (cf.~Algebra, Chapter VII, § 3, Exercise 8 (b)).

\begin{enumerate}
\def\labelenumi{(\alph{enumi})}
\setcounter{enumi}{1}
\tightlist
\item
  From (a) deduce an example of an exact sequence
\end{enumerate}

\[
0 \rightarrow \mathrm{E} ^ {\prime} \rightarrow \mathrm{E} \rightarrow \mathrm{E} ^ {\prime \prime} \rightarrow 0
\]

which does not split and whose terms are right A-modules which are not flat, such that for every left A-module F the sequence

\[
0 \rightarrow \mathrm{E} ^ {\prime} \otimes \mathrm{F} \rightarrow \mathrm{E} \otimes \mathrm{F} \rightarrow \mathrm{E} ^ {\prime \prime} \otimes \mathrm{F} \rightarrow 0
\]

is exact (use Lemma 2 of no. 1).

\begin{enumerate}
\def\labelenumi{\arabic{enumi}.}
\setcounter{enumi}{4}
\tightlist
\item
  Give an example of a right A-module E, a left A-module F and two submodules \(F'\) , \(F''\) of F such that the canonical image of \(E \otimes (F' \cap F'')\) in \(E \otimes F\) is not the intersection of the canonical images of \(E \otimes F'\) and \(E \otimes F''\) (cf.~Exercise 3(a)).
\end{enumerate}

¶ 6. Let A be a ring and M a left A-module. An exact sequence

\[
\mathrm{L} _ {n} \rightarrow \mathrm{L} _ {n - 1} \rightarrow \dots \rightarrow \mathrm{L} _ {1} \rightarrow \mathrm{L} _ {0} \rightarrow \mathrm{M} \rightarrow 0
\]

where \(L_{i}\) is a free left A-module ( \(0 \leqslant i \leqslant n\) ), is called a presentation of M of length n or an n-presentation of M. The presentation is called finite if all the \(L_{i}\) are finitely generated free modules.

If M is a finitely generated left A-module, we denote by \(\lambda(M)\) the least upper bound (finite or \(+\infty\) ) of the integers \(n \geqslant 0\) such that M has a finite \(n\) -presentation. If M is not finitely generated, we set \(\lambda(M) = -1\) .

\begin{enumerate}
\def\labelenumi{(\alph{enumi})}
\item
  Let \(0 \to \mathbf{P} \to \mathbf{N} \to \mathbf{M} \to 0\) be an exact sequence of left A-modules. Then \(\lambda(\mathbf{N}) \geqslant \inf(\lambda(\mathbf{P}), \lambda(\mathbf{M}))\) . (Starting with two \(n\) -presentations of \(\mathbf{P}\) and \(\mathbf{M}\) respectively, derive one for \(\mathbf{N}\) using Exercise 5 of § 1.)
\item
  Let \(M_{n} \xrightarrow{u_{n}} M_{n-1} \rightarrow \cdots \rightarrow M_{0} \xrightarrow{u_{0}} M \rightarrow 0\) be a finite n-presentation of M; show that, if \(\lambda(M) > n\) , \(\operatorname{Ker}(u_{n})\) is a finitely generated A-module. (Let
\end{enumerate}

\[
\mathrm{L} _ {n + 1} \xrightarrow {v _ {n + 1}} \mathrm{L} _ {n} \longrightarrow \dots \longrightarrow \mathrm{L} _ {2} \xrightarrow {v _ {2}} \mathrm{L} _ {1} \xrightarrow {v _ {1}} \mathrm{L} _ {0} \xrightarrow {v _ {0}} \mathrm{M} \longrightarrow 0
\]

be a finite \((n + 1)\) -presentation of \(\mathbf{M}\) and let \(\mathbf{P} = \operatorname{Ker}(v_0)\) be such that there is an \(n\) -presentation of \(\mathbf{P}\) :

\[
\mathrm{L} _ {n + 1} \xrightarrow {v _ {n + 1}} \mathrm{L} _ {n} \longrightarrow \dots \longrightarrow \mathrm{L} _ {2} \xrightarrow {v _ {2}} \mathrm{L} _ {1} \xrightarrow {v _ {1}} \mathrm{P} \longrightarrow 0.
\]

Applying the method of (a) to the exact sequence \(0 \rightarrow P \rightarrow L_{0} \rightarrow M \rightarrow 0\) we obtain an exact sequence

\[
\mathrm{M} _ {n} \oplus \mathrm{L} _ {n + 1} \xrightarrow {w _ {n}} \mathrm{M} _ {n - 1} \oplus \mathrm{L} _ {n} \xrightarrow {w _ {n - 1}} \dots \xrightarrow {w _ {1}} \mathrm{M} _ {0} \oplus \mathrm{L} _ {1} \xrightarrow {w _ {0}} \mathrm{L} _ {0} \longrightarrow 0
\]

and exact sequences \(0 \to \operatorname{Ker}(v_{i+1}) \to \operatorname{Ker}(w_i) \to \operatorname{Ker}(u_i) \to 0\) (§ 1, no. 4, Proposition 2). Observe finally that \(\operatorname{Ker}(w_i)\) is a direct factor of \(\mathbf{M}_i \oplus \mathbf{L}_{i+1}\) .

\begin{enumerate}
\def\labelenumi{(\alph{enumi})}
\setcounter{enumi}{2}
\tightlist
\item
  Show that with the hypotheses of (a)
\end{enumerate}

\[
\lambda (\mathbf {M}) \geqslant \inf (\lambda (\mathbf {N}), \lambda (\mathbf {P}) + 1).
\]

(If \(n \leqslant \inf(\lambda(N), \lambda(P) + 1)\) , show by induction on \(n\) that \(\lambda(M) \geqslant n\) , arguing as in (a) and using (b).)

\begin{enumerate}
\def\labelenumi{(\alph{enumi})}
\setcounter{enumi}{3}
\tightlist
\item
  Show that with the hypotheses of (a)
\end{enumerate}

\[
\lambda (\mathrm{P}) \geqslant \inf (\lambda (\mathrm{N}), \lambda (\mathrm{M}) - 1).
\]

(Same method as in (c).) Deduce that, if \(\lambda(N) = +\infty\) , then \(\lambda(M) = \lambda(P) + 1\) .

\begin{enumerate}
\def\labelenumi{(\alph{enumi})}
\setcounter{enumi}{4}
\tightlist
\item
  Deduce from (a), (c) and (d) that, if \(\mathbf{N} = \mathbf{M} \oplus \mathbf{P}\) , then
\end{enumerate}

\[
\lambda (\mathbf {N}) = \inf (\lambda (\mathbf {M}), \lambda (\mathbf {P})).
\]

In particular, for N to admit a finite presentation it is necessary and sufficient that M and P do so too.

\begin{enumerate}
\def\labelenumi{(\alph{enumi})}
\setcounter{enumi}{5}
\tightlist
\item
  Let \(N_{1}\) , \(N_{2}\) be two submodules of an A-module M. Suppose that \(N_{1}\) and \(N_{2}\) admit finite presentations. For \(N_{1} + N_{2}\) to admit a finite presentation it is necessary and sufficient that \(N_{1} \cap N_{2}\) be finitely generated.
\end{enumerate}

\begin{enumerate}
\def\labelenumi{\arabic{enumi}.}
\setcounter{enumi}{6}
\tightlist
\item
  \begin{enumerate}
  \def\labelenumii{(\alph{enumii})}
  \tightlist
  \item
    With the notation of Exercise 6, show that, if M is a projective module, then \(\lambda(M) = -1\) or \(\lambda(M) = +\infty\) . If A is a left Noetherian ring, then, for every A-module M, \(\lambda(M) = -1\) or \(\lambda(M) = +\infty\) .
  \end{enumerate}
\end{enumerate}

\begin{enumerate}
\def\labelenumi{(\alph{enumi})}
\setcounter{enumi}{1}
\item
  If \(\mathfrak{a}\) is a left ideal of a ring \(\mathbf{A}\) , which is not finitely generated, \(\mathbf{A}_s / \mathfrak{a}\) is a monogenous A-module which does not admit a finite presentation, in other words \(\lambda(\mathbf{A}_s / \mathfrak{a}) = 0\) (no. 8, Lemma 9).
\item
  Give an example of a monogenous left ideal \(\mathfrak{a}\) of a ring A such that \(\mathbf{A}_s / \mathfrak{a}\) (which has a finite presentation) admits a dual which is not a finitely generated right A-module.
\item
  Let K be a commutative field, E the vector space \(\mathbf{K}^{(\mathbf{N})}, (e_n)\) the canonical basis of E and T the tensor algebra of E, which therefore has a basis consisting of the finite products \(e_{i_1}e_{i_2}\ldots e_{i_k}\) ( \(k \geqslant 0, i_j \in N\) for all j). For a given integer n, let b be the two-sided ideal of T generated by the products
\end{enumerate}

\[
e _ {1} e _ {0}, e _ {2} e _ {1}, \dots , e _ {n} e _ {n - 1}
\]

and \(e_{n+k}e_n\) for all \(k\geqslant 1\) ; let A be the quotient ring T/b, and for every integer \(m\) , let \(a_m\) be the canonical image of \(e_m\) in A. Show that, if \(\mathbf{M} = \mathbf{A}_s / \mathbf{A}a_0\) , then \(\lambda (\mathbf{M}) = n\) (observe that, for \(m\leqslant n - 1\) , the left annihilator of \(a_m\) is \(\mathrm{A}a_{m+1}\) and use Exercise 6(b)).

\begin{enumerate}
\def\labelenumi{\arabic{enumi}.}
\setcounter{enumi}{7}
\item
  Let C be a commutative ring and E, F two C-modules. Show that \(\lambda(E \otimes_{\mathbb{C}} F) \geqslant \inf(\lambda(E), \lambda(F))\) .
\item
  Let E be a finitely presented left A-module.
\end{enumerate}

\begin{enumerate}
\def\labelenumi{(\alph{enumi})}
\item
  Show that for every family \((\mathbf{F}_t)_{t \in I}\) of right A-modules, the canonical homomorphism \(\mathbf{E} \otimes_{\mathbf{A}} \left( \prod_{t \in I} \mathbf{F}_t \right) \to \prod_{t \in I} (\mathbf{E} \otimes_{\mathbf{A}} \mathbf{F}_t)\) (Algebra, Chapter II, § 3, no. 7) is bijective.
\item
  Let \((G_{\alpha}, \phi_{\beta\alpha})\) be a direct system of left A-modules; show that the canonical homomorphism
\end{enumerate}

\[
\varinjlim \operatorname{Hom} _ {\mathrm{A}} (\mathrm{E}, \mathrm{G} _ {\alpha}) \rightarrow \operatorname{Hom} _ {\mathrm{A}} (\mathrm{E}, \varinjlim \mathrm{G} _ {\alpha})
\]

is bijective.

\begin{enumerate}
\def\labelenumi{\arabic{enumi}.}
\setcounter{enumi}{9}
\tightlist
\item
  \begin{enumerate}
  \def\labelenumii{(\alph{enumii})}
  \tightlist
  \item
    Let A be a ring, I a set and R a submodule of \(L = A_{d}^{(I)}\) . Let S be the set of ordered pairs (J, S), where J is a finite subset of I and S a finitely generated submodule of \(A_{d}^{J} \cap R\) . S is ordered by the relation ``J ⊂ J' and S ⊂ S'\,''; show that S is directed with respect to this order relation, that the family ( \(A_{d}^{J}/S\) ) is a direct system of right A-modules with S as indexing set and that there exists an isomorphism of L/R onto \(\varinjlim (A_{d}^{J}/S)\) .
  \end{enumerate}
\end{enumerate}

\begin{enumerate}
\def\labelenumi{(\alph{enumi})}
\setcounter{enumi}{1}
\tightlist
\item
  Deduce from (a) that every A-module is a direct limit of finitely presented A-modules.
\end{enumerate}

\begin{enumerate}
\def\labelenumi{\arabic{enumi}.}
\setcounter{enumi}{10}
\tightlist
\item
  Let E be a right A-module. E is called pseudo-coherent if every finitely generated submodule of E is finitely presented; every submodule of a pseudo-coherent module is pseudo-coherent. E is called coherent if it is pseudo-coherent and finitely generated (and therefore finitely presented).
\end{enumerate}

\begin{enumerate}
\def\labelenumi{(\alph{enumi})}
\item
  Let \(0 \to E' \to E \to E'' \to 0\) be an exact sequence of right A-modules. Show that, if E is pseudo-coherent (resp. coherent) and \(E'\) is finitely generated \(E''\) is pseudo-coherent (resp. coherent). Show that, if \(E'\) and \(E''\) are pseudo-coherent (resp. coherent), so is E. Show that, if E and \(E''\) are coherent, so is \(E'\) (use Exercise 6 and Lemma 9 of no. 8).
\item
  Let E be a coherent A-module and \(E'\) a pseudo-coherent (resp. coherent) A-module. Show that, for every homomorphism \(u: E \to E'\) , \(\text{Im}(u)\) and \(\text{Ker}(u)\) are coherent and that \(\text{Coker}(u)\) is pseudo-coherent (resp. coherent) (use (a)).
\item
  Show that every direct sum (resp. every finite direct sum) of pseudo-coherent (resp. coherent) modules is a pseudo-coherent (resp. coherent) module.
\item
  If E is a pseudo-coherent module and M, N are coherent submodules of E, show that \(M + N\) and \(M \cap N\) are coherent (use (a) and (c)).
\item
  Suppose that A is commutative. Show that, if E is a coherent A-module and F a coherent (resp. pseudo-coherent) A-module, \(\operatorname{Hom}_{\mathbf{A}}(\mathbf{E}, \mathbf{F})\) is a coherent (resp. pseudo-coherent) A-module. (Reduce it to the case where F is coherent and consider a finite presentation of E, then use (b).)
\end{enumerate}

¶ 12. (a) Let A be a ring. Show that the following four properties are equivalent:

\((\alpha)\) The right A-module \(\mathbf{A}_d\) is coherent (Exercise 11).

(β) Every finitely presented right A-module is coherent.

(γ) Every A-module \(A_{s}^{I}\) (I an arbitrary set) is flat.

(δ) Every product of flat left A-modules is flat.

(To prove that \((\alpha)\) implies \((\beta)\) , use Exercise 11(b). To see that \((\gamma)\) implies \((\alpha)\) , use Proposition 13 of no. 11 and argue by reductio ad absurdum. To show that \((\alpha)\) implies \((\delta)\) , use Exercises 9.)

Such a ring is called right coherent and the concept of a left coherent ring is defined similarly.

\begin{enumerate}
\def\labelenumi{(\alph{enumi})}
\setcounter{enumi}{1}
\tightlist
\item
  Show that every right Noetherian ring is right coherent. Give an example of a right Artinian ring which is not left coherent (cf.~Algebra, Chapter VIII, § 2, Exercise 4).
\end{enumerate}

*(c) The ring of a non-discrete valuation of height 1 (Chapter VI) is coherent but contains ideals which are not coherent and admits (monogenous) quotient modules which are not pseudo-coherent.*

\begin{enumerate}
\def\labelenumi{(\alph{enumi})}
\setcounter{enumi}{3}
\item
  Show that, if A is a right coherent ring, then, for every right A-module E, \(\lambda(\mathrm{E}) = -1\) or \(\lambda(\mathrm{E}) = 0\) or \(\lambda(\mathrm{E}) = +\infty\) (Exercise 6).
\item
  Let \((\mathbf{A}_{\alpha}, \phi_{\beta \alpha})\) be a direct system of rings whose indexing set is directed and let \(\mathbf{A} = \varinjlim \mathbf{A}_{\alpha}\) . Suppose that, for \(\alpha \leqslant \beta\) , \(\mathbf{A}_{\beta}\) is a flat right \(\mathbf{A}_{\alpha}\) -module. Show that, if the \(\mathbf{A}_{\alpha}\) are right coherent, so is \(\mathbf{A}\) . (Observe that \(\mathbf{A}\) is a flat \(\mathbf{A}_{\alpha}\) -module for all \(\alpha\) and that, if \(\mathbf{E}\) is a finitely generated submodule of \(\mathbf{A}_d\) , there exists an index \(\alpha\) and a finitely generated submodule \(\mathbf{E}_{\alpha}\) of \((\mathbf{A}_{\alpha})_d\) such that \(\mathbf{E}_{\alpha} \otimes_{\mathbf{A}_{\alpha}} \mathbf{A}\) is isomorphic to \(\mathbf{E}\) .)
\end{enumerate}

* (f) Deduce from (e) that every polynomial ring (in any finite or infinite set of indeterminates) over a Noetherian commutative ring is coherent. Deduce from this that a quotient ring of a coherent ring is not necessarily coherent.*

\begin{enumerate}
\def\labelenumi{(\alph{enumi})}
\setcounter{enumi}{6}
\tightlist
\item
  For A to be left coherent, it is necessary and sufficient that the left annihilator of every element of A be finitely generated and that the intersection of two finitely generated left ideals of A be finitely generated (use Exercise 6(f)).
\end{enumerate}

¶ 13. Let A, B be two rings, F an (A, B)-bimodule and G a right B-module. Show that, if G is injective (Algebra Chapter II, § 2, Exercise 11) and F is a flat left A-module, the right A-module \(\operatorname{Hom}_{B}(F, G)\) is injective. (Use the isomorphism

\[
\operatorname{Hom} _ {\mathbf {A}} (\mathbf {E}, \operatorname{Hom} _ {\mathbf {B}} (\mathbf {F}, \mathbf {G})) \rightarrow \operatorname{Hom} _ {\mathbf {B}} (\mathbf {E} \otimes_ {\mathbf {A}} \mathbf {F}, \mathbf {G})
\]

for a right A-module E (Algebra, Chapter II, § 4, no. 1).)

¶ 14. Let A, B be two rings, E a left A-module, F an (A, B)-bimodule and G a right B-module; consider the canonical homomorphism (Algebra, Chapter II, § 4, Exercise 5)

\[
\sigma \colon \operatorname{Hom} _ {\mathbb {B}} (\mathrm{F}, \mathrm{G}) \otimes_ {\mathrm{A}} \mathrm{E} \rightarrow \operatorname{Hom} _ {\mathbb {B}} (\operatorname{Hom} _ {\mathrm{A}} (\mathrm{E}, \mathrm{F}), \mathrm{G})
\]

such that \((\sigma(u \otimes x))(v) = u(v(x))\) for all \(x \in \mathbf{E}\) , \(u \in \operatorname{Hom}_{\mathbf{B}}(\mathbf{F}, \mathbf{G})\) ,

\[
v \in \operatorname{Hom} _ {\mathbf {A}} (\mathrm{E}, \mathbf {F}).
\]

Show that, if G is an injective B-module (Algebra, Chapter II, § 2, Exercise 11) and E finitely presented, \(\sigma\) is bijective. (Consider first the case where E is free and finitely generated.)

¶ 15. Let A be a ring. Show that every left A-module E which is flat and finitely presented is projective. (Given a surjective left A-module homomorphism \(u: \mathbf{F} \to \mathbf{F}'\) , let \(u'\) be the homomorphism

\[
\operatorname{Hom} \left(1 _ {\mathbf {E}}, u\right): \operatorname{Hom} _ {\mathbf {A}} (\mathrm{E}, \mathrm{F}) \rightarrow \operatorname{Hom} _ {\mathbf {A}} (\mathrm{E}, \mathrm{F} ^ {\prime \prime})
\]

and \(\bar{u}\) the homomorphism

\[
\operatorname{Hom} \left(u ^ {\prime}, 1 _ {G}\right): \operatorname{Hom} _ {\mathbf {Z}} \left(\operatorname{Hom} _ {A} (E, F ^ {\prime \prime}), G\right)\rightarrow \operatorname{Hom} _ {\mathbf {Z}} \left(\operatorname{Hom} _ {A} (E, F), G\right),
\]

where G is a divisible Z-module. Using Exercise 14 first, prove that \(\bar{u}\) is injective; then, with a suitable choice for G (Algebra, Chapter II, §2, Exercise 14), show that \(u'\) is surjective.)

\begin{enumerate}
\def\labelenumi{\arabic{enumi}.}
\setcounter{enumi}{15}
\tightlist
\item
  Let A be a ring and \(a\) an element of A. Show that the following properties are equivalent.
\end{enumerate}

\((\alpha)\) \(a \in aAa\) .

(β) aA is a direct factor of the module \(A_{d}\) .

(γ) \(A_{d}/aA\) is a flat right A-module.

(δ) For every left ideal b of A, \(aA \cap b = ab\) .

(Use the Corollary to Proposition 7 of no. 6 to prove the equivalence of \((\gamma)\) and \((\delta)\) and show directly that \((\delta)\) implies \((\alpha)\) and \((\alpha)\) implies \((\beta)\) , by proving the existence of an idempotent element \(e \in a\mathbf{A}\) such that \(e\mathbf{A} = a\mathbf{A}\) .)

\begin{enumerate}
\def\labelenumi{\arabic{enumi}.}
\setcounter{enumi}{16}
\tightlist
\item
  Let A be a ring. Show that the following properties are equivalent:
\end{enumerate}

(α) Every element \(a \in A\) satisfies the equivalent properties of Exercise 16.

(β) Every finitely generated right ideal of A is a direct factor of \(A_{d}\) .

(γ) Every left A-module is flat.

(δ) Every right A-module is flat.

Then A is called an absolutely flat ring (*).

(To see that \((\alpha)\) implies \((\beta)\) , use Exercise 15(b) of Algebra, Chapter VIII, §6.)

¶ 18. Let A be an absolutely flat ring (Exercise 17).

\begin{enumerate}
\def\labelenumi{(\alph{enumi})}
\item
  Let P be a projective right A-module. Show that every finitely generated submodule E of P is a direct factor of P. (Reduce it to the case where P is free and finitely generated. Then note that P/E is finitely presented and use Exercises 15.)
\item
  Show that every projective right A-module P is the direct sum of monogenous submodules isomorphic to monogenous right ideals of A. (Use Kaplansky's Theorem (Algebra, Chapter II, § 2, Exercises 3) to reduce the problem to the case where P is generated by a countable family of elements, then use (a).)
\item
  Give an example of an absolutely flat ring A and a non-projective finitely generated A-module (consider a quotient of A by an ideal which is not finitely generated; cf.~Algebra, Chapter VIII, § 6, Exercise 15(f) and Commutative Algebra, Chapter II, § 4, Exercise 17).
\end{enumerate}

\begin{enumerate}
\def\labelenumi{\arabic{enumi}.}
\setcounter{enumi}{18}
\tightlist
\item
  Let A be a ring. Show that the following properties are equivalent:
\end{enumerate}

\((\alpha)\) A is semisimple.

(β) Every right ideal of A is an injective A-module.

(γ) Every right A-module is projective.

(δ) Every right A-module is injective.

\begin{enumerate}
\def\labelenumi{\arabic{enumi}.}
\setcounter{enumi}{19}
\item
  Let A be an integral domain, B an A-algebra which is a flat A-module and M a torsion-free A-module. Show that, if \(t \in B\) is not a divisor of zero, the relation t. z = 0 for \(z \in B \otimes_{A} M\) implies z = 0. (Reduce this to the case where M is finitely generated and, by embedding M in a finitely generated free A-module, to the case where M = A.)
\item
  Let S be a commutative ring, R a commutative S-algebra, B an S-algebra (commutative or otherwise) and \(\mathrm{B}_{(\mathrm{R})}\) the R-algebra obtained from B by extension of scalars. Suppose that R is a flat S-module and that B is a finitely generated S-module. If Z is the centre of B, show that the canonical homomorphism of \(Z_{(R)} = Z \otimes_{S} R\) to \(\mathrm{B}_{(\mathrm{R})}\) is an isomorphism of \(Z_{(R)}\) onto the centre of \(\mathrm{B}_{(\mathrm{R})}\) . (Use the exact sequence
\end{enumerate}

\[
0 \longrightarrow \mathrm{Z} \longrightarrow \mathrm{B} \stackrel {\theta} {\longrightarrow} \operatorname{Hom} _ {\mathrm{s}} (\mathrm{B}, \mathrm{B})
\]

where \(\theta (x)(y) = xy - yx\) , and Proposition 11 of no. 10.)

\begin{enumerate}
\def\labelenumi{\arabic{enumi}.}
\setcounter{enumi}{21}
\tightlist
\item
  Let E be a left A-module. For every right ideal a of A and every element \(a \in A\) , denote by a: a the set of \(x \in A\) such that \(ax \in a\) and by aE: a the set of \(y \in E\) such that \(ay \in aE\) . Then clearly (a: a)E ⊂ aE: a. Show that, for E to be flat, it is necessary and sufficient that, for every right ideal a of A and every element \(a \in A\) , (a: a)E = aE: a. (To see that the condition is necessary, consider the exact sequence of right A-modules \(0 \to (a: a)/a \stackrel{\psi}{\to} A_d/a \stackrel{\varphi}{\to} A_d/a\) , where \(\psi\) is the canonical injection and \(\phi\) the mapping obtained by taking quotients under left multiplication by a. To see that the condition is sufficient apply the criterion of Corollary 1 to Proposition 13 of no. 11: starting with a relation \(\sum_{i=1}^{n} a_i x_i = 0\) , where \(a_i \in A\) , \(x_i \in E\) , apply the hypothesis to the ideal
\end{enumerate}

\(a_{2}=\sum_{i=2}^{n}a_{i}A\) and the element \(a_{1}\) and argue by induction on n.)

\begin{enumerate}
\def\labelenumi{\arabic{enumi}.}
\setcounter{enumi}{22}
\tightlist
\item
  \begin{enumerate}
  \def\labelenumii{(\alph{enumii})}
  \tightlist
  \item
    Let \(0 \to \mathbf{R} \to \mathbf{L} \to \mathbf{E} \to 0\) be an exact sequence of left A-modules, where \(\mathbf{L}\) is a free A-module; let \((e_{\alpha})\) be a basis of \(\mathbf{L}\) . Show that the following conditions are equivalent:
  \end{enumerate}
\end{enumerate}

(α) E is flat.

(β) For all \(x \in \mathbb{R}\) , if \(\mathfrak{a}_x\) is the right ideal generated by the components of \(x\) with respect to the basis \((e_{\alpha})\) , then \(x \in \mathbb{R}\mathfrak{a}_x\) .

\((\gamma)\) For all \(x \in \mathbb{R}\) , there exists a homomorphism \(u_x: \mathbf{L} \to \mathbf{R}\) such that \(u_x(x) = x\) .

\begin{enumerate}
\def\labelenumi{(\arabic{enumi})}
\setcounter{enumi}{7}
\tightlist
\item
  For every finite sequence \((x_{i})_{1\leqslant i\leqslant n}\) of elements of R, there exists a homomorphism \(u\colon \mathbf{L}\to \mathbb{R}\) such that \(u(x_{i}) = x_{i}\) for \(1\leqslant i\leqslant n\) . (Use the Corollary to Proposition 7 of no. 6.)
\end{enumerate}

\begin{enumerate}
\def\labelenumi{(\alph{enumi})}
\setcounter{enumi}{1}
\tightlist
\item
  Let \(\mathfrak{a}\) be a left ideal of \(\mathbf{A}\) such that \(\mathbf{A} / \mathfrak{a}\) is a flat \(\mathbf{A}\) -module. Show that, for every finitely generated left ideal \(\mathfrak{b} \subset \mathfrak{a}\) , there exists \(x \in \mathbf{A}\) such that
\end{enumerate}

\[
\mathfrak {b} \subset \mathrm{A} x \subset \mathfrak {a}
\]

(use condition (δ) of (a)).

\begin{enumerate}
\def\labelenumi{(\alph{enumi})}
\setcounter{enumi}{2}
\item
  Derive from (a) a new proof of the result of Exercise 15.
\item
  Let \(\mathfrak{r}\) be the Jacobson radical of \(\mathbf{A}\) and \(0 \to \mathbf{R} \to \mathbf{L} \to \mathbf{E} \to 0\) an exact sequence of left A-modules such that \(\mathbf{L}\) is free. Suppose that \(\mathbf{E}\) is flat and \(\mathbf{R}\) is contained in \(\mathfrak{rL}\) . Show that \(\mathbf{R} = 0\) (in the notation of (a) observe that \(\mathfrak{a}_x\) is a finitely generated ideal and that \(\mathfrak{a}_x = \mathfrak{a}_x\mathfrak{r}\) ).
\item
  Let E be a finitely generated flat A-module; suppose that there exists a two-sided ideal b of A contained in the Jacobson radical of A such that E/bE is a free (A/b)-module. Show that E is then a free A-module (observe that there exists a finitely generated free A-module L such that L/bL is isomorphic to E/bE and use Algebra, Chapter VIII, § 6, no. 3, Corollary 4 to Proposition 6; then apply (d)).
\end{enumerate}

\begin{enumerate}
\def\labelenumi{\arabic{enumi}.}
\setcounter{enumi}{23}
\tightlist
\item
  A submodule \(M'\) of a right A-module M is called pure if, denoting by \(j: M' \rightarrow M\) the canonical injection, the homomorphism
\end{enumerate}

\[
j \otimes 1 _ {N} \colon M ^ {\prime} \otimes_ {A} N \rightarrow M \otimes_ {A} N
\]

is injective for every left A-module N. This is so if \(M'\) is a direct factor of M or if \(M/M'\) is flat, but these two conditions are not necessary (Exercise 4).

\begin{enumerate}
\def\labelenumi{(\alph{enumi})}
\item
  Show that, for \(\mathbf{M}'\) to be a pure submodule of \(\mathbf{M}\) , it is necessary and sufficient that, if \((m_i')_{i\in I}\) is a finite family of elements of \(\mathbf{M}'\) , \((x_j)_{j\in J}\) a family of elements of \(\mathbf{M}\) such that \(m_i' = \sum_{j\in J} x_j a_{jl}\) for all \(i\in I\) and a family \((a_{jl})\) of elements of \(\mathbf{A}\) , then there exists a family \((x_j')_{j\in J}\) of elements of \(\mathbf{M}'\) such that \(m_i' = \sum_{j\in J} x_j' a_{jl}\) for all \(i\in I\) . (To see that the condition is sufficient, use Lemma 10 of no. 11 to show that \(\mathbf{M}' \otimes_{\mathbf{A}} \mathbf{N} \to \mathbf{M} \otimes_{\mathbf{A}} \mathbf{N}\) is injective for every finitely generated left A-module \(\mathbf{N}\) ; to see that the condition is necessary, consider a finitely generated left A-module \(\mathbf{N} = \mathbf{L}/\mathbf{R}\) , where \(\mathbf{L}\) is a finitely generated free A-module and \(\mathbf{R}\) is a finitely generated submodule of \(\mathbf{L}\) .) Deduce from this criterion that, if \(\mathbf{A}\) is a principal ideal domain, the notion of pure submodule of an A-module coincides with that of Algebra, Chapter VII, § 2, Exercise 7.
\item
  Let M be a right A-module, M' a submodule of M and M'' a submodule of M'. Show that, if M' is a pure submodule of M and M'' a pure submodule of M', then M'' is a pure submodule of M and M'/M'' is a pure submodule of M/M''. If M'' is a pure submodule of M, M'' is a pure submodule of M'.
\item
  Show that, if N and P are two submodules of M such that \(N \cap P\) and \(N + P\) are pure in M, then N and P are pure submodules of M. Give an example of two submodules N, P of \(Z^{2}\) which are pure in \(Z^{2}\) but where \(N + P\) is not a pure submodule of \(Z^{2}\) .
\item
  Let C be a commutative ring and E, F two C-modules; show that, if \(E'\) (resp. \(F'\) ) is a pure submodule of E (resp. F), the canonical mapping \(E' \otimes_{C} F' \to E \otimes_{C} F\) is injective and identifies \(E' \otimes_{C} F'\) with a pure submodule of \(E \otimes_{C} F\) .
\item
  Let \(\rho: \mathbf{A} \to \mathbf{B}\) be a ring homomorphism, \(\mathbf{M}\) a right A-module and \(\mathbf{M}'\) a pure submodule of \(\mathbf{M}\) . Show that \(\mathbf{M}_{(\mathbf{B})}' = \mathbf{M}' \otimes_{\mathbf{A}} \mathbf{B}\) is canonically identified with a pure submodule of \(\mathbf{M}_{(\mathbf{B})} = \mathbf{M} \otimes_{\mathbf{A}} \mathbf{B}\) .
\end{enumerate}

